\section{Conclusões}

Este trabalho apresentou uma abordagem integrada de AM e PO para apoiar o processo de montagem de elencos no futebol profissional. 
Inicialmente, foi desenvolvido um modelo preditivo capaz de estimar o desempenho esperado de equipes a partir das características individuais dos jogadores. 
Os resultados obtidos indicam que o modelo proposto alcança desempenho aproximadamente 55\% superior ao baseline adotado, evidenciando sua capacidade de capturar 
padrões relevantes nos dados e produzir estimativas mais precisas de desempenho coletivo.

A partir dessas previsões, o modelo foi utilizado como função objetivo em um problema de otimização voltado à seleção de jogadores sob restrições orçamentárias. 
Para explorar o espaço de soluções, foram implementadas duas heurísticas: um algoritmo guloso baseado em substituições iterativas e um método de busca local.
\begin{comment}
Os experimentos realizados indicaram que ambos os algoritmos são capazes de produzir soluções de alta qualidade, com diferenças muito pequenas nas pontuações 
previstas obtidas. Embora a busca local explore o espaço de soluções de forma mais abrangente, os resultados mostram que a heurística gulosa é capaz de alcançar 
soluções comparáveis em termos de desempenho previsto. Além disso, o algoritmo guloso apresentou comportamento ligeiramente mais estável entre execuções e 
realizou, em média, um número semelhante ou menor de avaliações da função objetivo. Esses resultados sugerem que, para o problema analisado, o espaço de soluções 
apresenta uma estrutura relativamente suave, na qual heurísticas simples já são capazes de identificar soluções próximas às obtidas por métodos de busca mais 
elaborados.  
\end{comment}
\textit{Inserir resultados dos algoritmos de PO}

Do ponto de vista prático, os resultados obtidos indicam que a abordagem proposta pode contribuir para tornar o processo de tomada de decisão na montagem de 
elencos mais estruturado e orientado por dados. No contexto do futebol brasileiro, onde decisões de contratação frequentemente são influenciadas por fatores 
subjetivos ou avaliações intuitivas, ferramentas baseadas em análise quantitativa podem oferecer suporte adicional à identificação de jogadores que se encaixem 
de forma mais eficiente em determinadas estruturas de equipe, respeitando restrições financeiras.

Assim, este trabalho demonstra que a integração entre técnicas de AM e métodos de otimização representa um caminho promissor para apoiar 
processos de recrutamento e planejamento esportivo. Ao transformar dados de desempenho em estimativas de impacto coletivo e utilizá-las em modelos de decisão, 
torna-se possível reduzir a dependência exclusiva da intuição na escolha de atletas e avançar em direção a abordagens mais analíticas na gestão do futebol.